\chapter{Tình Người Trong Gian Khó (Humanity in Hardship)}
\begin{center}
  \textit{\color{truyengold}``Vàng thật không sợ lửa --- người thật không sợ gian khó.''}\\[4pt]
  \textit{(True gold does not fear fire --- a true person does not fear hardship.)}\\[4pt]
  \small\color{truyenblue}--- Tục ngữ Việt Nam ---
\end{center}
\ngancach

\section{Bát Cháo Hành Của Thị Nở}
\begin{truyen}{Bát Cháo Hành Của Thị Nở}{Nam Cao --- Việt Nam, 1941}
\chuhoa{C}{hí} Phèo --- kẻ bị xã hội ruồng bỏ, biến thành con quỷ làng Vũ Đại --- lần đầu tiên trong cuộc đời nhận được một bát cháo hành từ tay Thị Nở, người đàn bà xấu xí nhất làng.\\[4pt]
\textit{(Chi Pheo --- abandoned by society, turned into the village devil of Vu Dai --- received for the first time in his life a bowl of scallion porridge from Thi No, the plainest woman in the village.)}

Đó là lần đầu tiên trong đời hắn được người ta quan tâm --- không phải vì sợ, không phải vì thương hại, mà vì chân thật.\\[4pt]
\textit{(It was the first time in his life someone cared for him --- not out of fear, not out of pity, but out of genuine feeling.)}

Chí Phèo ngồi khóc vì bát cháo ấy, vì hắn cảm nhận lại được hơi ấm của con người sau bao nhiêu năm đắm chìm trong thú tính.\\[4pt]
\textit{(Chi Pheo sat and wept over that bowl, for he felt human warmth again after years drowning in brutality.)}

Nam Cao muốn nhắn nhủ: dù ai đó đã bị xã hội vứt bỏ đến đâu, một hành động yêu thương nhỏ bé cũng có thể đánh thức lại phần người trong họ.\\[4pt]
\textit{(Nam Cao wanted to say: no matter how far someone has been cast aside by society, a small act of love can still awaken the human in them.)}

\end{truyen}
\begin{baihoc}
  Không ai hoàn toàn mất đi phần người --- đôi khi chỉ cần một bát cháo hành để tìm lại chính mình.\\[4pt]
  \textit{(No one completely loses their humanity --- sometimes all it takes is a bowl of porridge to find oneself again.)}
\end{baihoc}

\section{Người Lính Nhật Bản Trao Nước Cho Tù Binh}
\begin{truyen}{Người Lính Nhật Bản Trao Nước Cho Tù Binh}{Thế Chiến II, 1944}
\chuhoa{T}{rong} đoàn tù binh Đồng Minh bị áp giải qua một ngôi làng Nhật Bản năm 1944, một người lính Nhật trẻ bí mật trao một bình nước cho người tù binh Mỹ kiệt sức.\\[4pt]
\textit{(Among Allied POWs being marched through a Japanese village in 1944, a young Japanese soldier secretly handed a flask of water to an exhausted American prisoner.)}

Hành động đó vi phạm lệnh; người lính có thể bị kỷ luật nếu bị phát hiện.\\[4pt]
\textit{(The act violated orders; the soldier could have been punished if discovered.)}

Người tù binh Mỹ ấy sống sót sau chiến tranh, và trong hồi ký viết năm 1990, ông dành cả một chương để kể lại khoảnh khắc đó: ``Tôi không biết tên anh. Nhưng anh đã dạy tôi rằng giữa hai chiến tuyến, vẫn có con người.''\\[4pt]
\textit{(That American prisoner survived the war, and in his 1990 memoir devoted an entire chapter to that moment: `I never knew his name. But he taught me that even between two battle lines, there is still humanity.')}

\end{truyen}
\begin{baihoc}
  Lòng nhân ái không chọn chiến tuyến --- nó chỉ nhìn thấy một con người đang cần giúp đỡ.\\[4pt]
  \textit{(Compassion does not choose sides --- it only sees a person who needs help.)}
\end{baihoc}

\section{Cây Tre Giữa Bão}
\begin{truyen}{Cây Tre Giữa Bão}{Thiên nhiên Việt Nam}
\chuhoa{K}{hi} bão lớn ập vào làng, mọi cây to cứng rễ đều bị bật gốc --- chỉ có bụi tre vẫn đứng.\\[4pt]
\textit{(When a great storm struck the village, every large rigid-rooted tree was uprooted --- only the bamboo clumps stood.)}

Tre không cứng mà dẻo; nó cúi xuống đến tận đất khi gió mạnh nhất, rồi khi gió qua đi, đứng thẳng lên như chưa từng ngã.\\[4pt]
\textit{(Bamboo is not rigid but flexible; it bends to the ground in the strongest winds, and when the wind passes, stands tall as though it never bowed.)}

Người dân Việt Nam nhìn tre và học: ``Dẻo dai mới bền --- cứng quá thì gãy.''\\[4pt]
\textit{(Vietnamese people look at bamboo and learn: `Flexibility endures --- rigidity breaks.')}

Qua hàng ngàn năm ngoại xâm, dân tộc Việt Nam sống sót không phải vì cứng rắn như đá --- mà vì dẻo dai như tre.\\[4pt]
\textit{(Through thousands of years of invasion, the Vietnamese people survived not by being hard as stone --- but by being resilient as bamboo.)}

\end{truyen}
\begin{baihoc}
  Mềm mại không phải yếu đuối --- đó là trí tuệ của sự tồn tại lâu dài.\\[4pt]
  \textit{(Flexibility is not weakness --- it is the wisdom of long survival.)}
\end{baihoc}

\section{Mẹ Teresa Và Người Hấp Hối Trên Đường Phố}
\begin{truyen}{Mẹ Teresa Và Người Hấp Hối Trên Đường Phố}{Mẹ Teresa --- Ấn Độ, 1952}
\chuhoa{N}{ăm} 1952, Mẹ Teresa nhặt một người đàn ông hấp hối nằm trên đường phố Calcutta --- cơ thể đầy giòi bọ, bị mọi người tránh xa.\\[4pt]
\textit{(In 1952, Mother Teresa found a dying man lying on a Calcutta street --- his body covered with worms, avoided by all.)}

Bà đưa ông vào nơi trú ẩn, rửa sạch vết thương, ngồi bên cạnh đến lúc ông trút hơi thở cuối cùng.\\[4pt]
\textit{(She brought him to shelter, cleaned his wounds, and sat beside him until he breathed his last.)}

Trước khi mất, người đàn ông nói nhỏ: ``Cả cuộc đời tôi sống như con vật --- nhưng tôi chết như một con người.''\\[4pt]
\textit{(Before he died, the man whispered: `All my life I lived like an animal --- but I am dying like a human being.')}

Mẹ Teresa nói: ``Điều tệ nhất không phải là đói khát hay nghèo khổ --- mà là cảm thấy mình không được ai muốn, không được ai yêu thương, không được ai quan tâm.''\\[4pt]
\textit{(Mother Teresa said: `The greatest evil is not hunger or poverty --- it is feeling unwanted, unloved, and uncared for.')}

\end{truyen}
\begin{baihoc}
  Phẩm giá của mỗi con người xứng đáng được tôn trọng --- dù họ là ai, dù họ đang ở đâu.\\[4pt]
  \textit{(The dignity of every human being deserves respect --- no matter who they are or where they are.)}
\end{baihoc}
